\subsection{序环}

定义 1.——给定交换环 A，我们称 A 上的一个序与 A 的环结构相容，如果它与 A 的加法群结构相容，并且满足以下公理：

(OR) 关系 \(x \geqslant 0\) 与 \(y \geqslant 0\) 蕴含 \(xy \geqslant 0\) 。

环 A 连同这样一个序称为序环。

例子。—— 1）环 Q 和 Z 赋予通常的序，是序环。2）序环的乘积，赋予积序，是序环。特别地，从集合 E 到序环 A 的映射全体所成的环 \(A^{E}\) 是一个序环。

\begin{enumerate}
\def\labelenumi{\arabic{enumi})}
\setcounter{enumi}{2}
\tightlist
\item
  序环的子环，赋予诱导序，也是序环。
\end{enumerate}

在序环中，关系 \(x \geqslant y\) 与 \(z \geqslant 0\) 蕴含 \(xz \geqslant yz\) 。事实上，这些不等式分别等价于 \(x - y \geqslant 0\) , \(z \geqslant 0\) 与 \((x - y)z \geqslant 0\) 。

类似地，可以证明关系 \(x \leqslant 0\) 与 \(y \geqslant 0\) （相应地 \(y \leqslant 0\) ）蕴含 \(xy \leqslant 0\) （相应地 \(xy \geqslant 0\) ）。这些结果常以符号法则的名义被引用（两个元素若同为 \(\geqslant 0\) 或同为 \(\leqslant 0\) ，则称它们同号）。由此推出，若 A 是一个全序环，则每个平方都是正的，特别地，每个幂等元（例如单位元）都是正的。

例。—— 在 Z 上只存在一种全序环结构：事实上 1 \textgreater{} 0，从而 n \textgreater{} 0 对每个自然数 \(n \neq 0\) 成立，由归纳法可得。相形之下，Z 上存在并非全序的序环结构（见下文）。

设 P 为序环 A 的正元素集。已知（VI，第 3 页，命题 3）P 决定 A 上的序关系。说 A 是一个序环，等价于说 P 满足以下性质：

\[
\begin{array}{l l} \left(\mathrm{AP} _ {\mathrm{I}}\right) & \mathrm{P} + \mathrm{P} \subset \mathrm{P} \\ \left(\mathrm{AP} _ {\mathrm{II}}\right) & \mathrm{PP} \subset \mathrm{P} \\ \left(\mathrm{AP} _ {\mathrm{III}}\right) & \mathrm{P} \cap (- \mathrm{P}) = \{0 \}. \end{array}
\]

确实，\((\mathrm{AP}_{\mathrm{I}})\) 与 \((\mathrm{AP}_{\mathrm{III}})\) 断言 A 的加法群是一个序群（VI，第 3 页，命题 3），而 \((\mathrm{AP}_{\mathrm{II}})\) 是 (OR) 的翻译。

回忆一下，A 上的序关系为全序的必要且充分条件是：

\[
\left(\mathrm{AP} _ {\mathrm{IV}}\right) \quad \mathrm{P} \cup (- \mathrm{P}) = \mathrm{A}.
\]

例. --- 在 Z 中，若取 P 为（通常意义下的）正偶数集合，则得到一个并非全序的环。

还要记得，在全序阿贝尔群中，关系 \(n \cdot x = 0\) （对自然数 \(n \neq 0\) ）蕴含 \(x = 0\) （VI，第 4 页）；由此我们得到以下结果。

命题 1. --- 全序环作为 Z-模是无挠的（II，第 313 页）。

